\documentclass[12pt,a4paper,twocolumn]{article}
\usepackage[utf8]{inputenc}
\usepackage[spanish]{babel}
\usepackage{amsmath}
\usepackage{amssymb}
\usepackage{physics}
\usepackage{geometry}
\geometry{margin=0.5cm}

\title{Formulario de Ecuaciones Teoría Electromagnética Tema 1}
\author{Emerson Warhman}
\date{}

\begin{document}
\small
\maketitle

\section*{Operaciones Vectoriales}

\subsection*{Coordenadas Polares}

\subsubsection*{Relación de Variables con Cartesianas}

De Esféricas a Cartesianas
\begin{align*}
    x & = r \sin \theta \cos \phi \\
    y & = r \sin \theta \sin \phi \\
    z & = r \cos \theta
\end{align*}

De Cartesianas a Esféricas
\begin{align*}
    r      & = \sqrt{x^2 + y^2 + z^2} \quad (r \geq 0)                                               \\
    \theta & = \cos^{-1} \frac{z}{\sqrt{x^2 + y^2 + z^2}} \quad (0^\circ \leq \theta \leq 180^\circ) \\
    \phi   & = \tan^{-1} \frac{y}{x}
\end{align*}

\subsection*{Coordenadas Cilíndricas}

\subsubsection*{Relación de Variables con Cartesianas}

De Cilíndricas a Cartesianas
\begin{align*}
    x & = \rho \cos \phi \\
    y & = \rho \sin \phi \\
    z & = z
\end{align*}

De Cartesianas a Cilíndricas
\begin{align*}
    \rho & = \sqrt{x^2 + y^2} \quad (\rho \geq 0) \\
    \phi & = \tan^{-1} \frac{y}{x}                \\
    z    & = z
\end{align*}

\subsection*{Elementos Diferenciales de Volumen y Área Superficial}

Para coordenadas esféricas:
\begin{equation*}
    dV = r^2 \sin \theta \, dr \, d\theta \, d\phi
\end{equation*}

\begin{itemize}
    \item Para una superficie esférica constante $r$:
          \begin{equation*}
              dA = r^2 \sin \theta \, d\theta \, d\phi
          \end{equation*}
    \item Para una superficie cónica constante $\theta$:
          \begin{equation*}
              dA = r \sin \theta \, dr \, d\phi
          \end{equation*}
    \item Para una superficie meridional constante $\phi$:
          \begin{equation*}
              dA = r \, dr \, d\theta
          \end{equation*}
\end{itemize}

Para coordenadas cilíndricas:
\begin{equation*}
    dV = \rho \, d\rho \, d\phi \, dz
\end{equation*}

\begin{itemize}
    \item Para una superficie de radio constante $\rho$:
          \begin{equation*}
              dA = \rho \, d\phi \, dz
          \end{equation*}
    \item Para una superficie angular constante $\phi$:
          \begin{equation*}
              dA = d\rho \, dz
          \end{equation*}
    \item Para una superficie axial constante $z$:
          \begin{equation*}
              dA = \rho \, d\rho \, d\phi
          \end{equation*}
\end{itemize}

\subsection*{Producto Punto}

\begin{equation*}
    \vec{A} \cdot \vec{B} = A_x B_x + A_y B_y + A_z B_z = AB \cos \theta
\end{equation*}

\subsubsection*{Proyección de un Vector sobre Otro}
\begin{equation*}
    \text{Proy}_{\vec{B}} \vec{A} = \frac{\vec{A} \cdot \vec{B}}{|\vec{B}|^2}
\end{equation*}

\subsection*{Producto Cruz}

\begin{equation*}
    \vec{A} \times \vec{B} = (A_y B_z - A_z B_y) \hat{a}_x + (A_z B_x - A_x B_z) \hat{a}_y + (A_x B_y - A_y B_x) \hat{a}_z
\end{equation*}

\subsection*{Productos Punto de Vectores Unitarios}

\section*{Operadores Diferenciales}

\subsection*{Operador Gradiente}

\begin{equation*}
    \nabla f = \frac{1}{h_1} \frac{\partial f}{\partial u_1} \hat{a}_1 + \frac{1}{h_2} \frac{\partial f}{\partial u_2} \hat{a}_2 + \frac{1}{h_3} \frac{\partial f}{\partial u_3} \hat{a}_3
\end{equation*}

\subsection*{Operador Divergencia}

\begin{equation*}
    \nabla \cdot \vec{F} = \frac{1}{h_1 h_2 h_3} \left[ \frac{\partial}{\partial u_1}(h_2 h_3 F_1) + \frac{\partial}{\partial u_2}(h_1 h_3 F_2) + \frac{\partial}{\partial u_3}(h_1 h_2 F_3) \right]
\end{equation*}

\subsection*{Operador Rotacional}

\begin{equation*}
    \nabla \times \vec{A} = \frac{1}{h_1 h_2 h_3}
    \begin{vmatrix}
        h_1 \hat{e}_1                 & h_2 \hat{e}_2                 & h_3 \hat{e}_3                 \\
        \frac{\partial}{\partial u_1} & \frac{\partial}{\partial u_2} & \frac{\partial}{\partial u_3} \\
        h_1 A_1                       & h_2 A_2                       & h_3 A_3
    \end{vmatrix}
\end{equation*}

\section*{Cálculo de Integral de Línea}

\subsection*{Integral de Línea de una Función Escalar}

\begin{equation*}
    \int_C f(x, y, z) \, ds
\end{equation*}

donde
\begin{equation*}
    ds = \sqrt{\left(\frac{dx}{dt}\right)^2 + \left(\frac{dy}{dt}\right)^2 + \left(\frac{dz}{dt}\right)^2} \, dt
\end{equation*}

\begin{equation*}
    \int_C f(x, y, z) \, ds = \int_a^b f(x(t), y(t), z(t)) \left| \frac{d\vec{r}}{dt} \right| dt
\end{equation*}

\subsection*{Integral de Línea de un Campo Vectorial}

\begin{equation*}
    \int_C \vec{F} \cdot d\vec{r} = \int_C F_x \, dx + F_y \, dy + F_z \, dz
\end{equation*}

donde
\begin{equation*}
    d\vec{r} = \frac{d\vec{r}}{dt} dt = \left( \frac{dx}{dt} \hat{i} + \frac{dy}{dt} \hat{j} + \frac{dz}{dt} \hat{k} \right) dt
\end{equation*}

\begin{equation*}
    \int_C \vec{F} \cdot d\vec{r} = \int_a^b \vec{F}(\vec{r}(t)) \cdot \frac{d\vec{r}}{dt} \, dt
\end{equation*}

\section*{Ley de Coulomb}

\subsection*{Densidades de Carga}

\begin{itemize}
    \item Densidad volumétrica $\rho$: $\rho = \frac{dq}{dV}$
    \item Densidad superficial $\sigma$: $\sigma = \frac{dq}{dA}$
    \item Densidad lineal $\lambda$: $\lambda = \frac{dq}{dl}$
\end{itemize}

\subsection*{Fuerza entre dos cargas}

\begin{equation*}
    \vec{F}_{12} = \frac{1}{4\pi \epsilon_0} \frac{q_1 q_2}{r_{12}^2} \hat{r}_{12}
\end{equation*}

\section*{Campo Eléctrico}

\subsection*{Campo Eléctrico debido a una Carga Discreta}

\begin{equation*}
    \vec{E} = \frac{1}{4\pi \epsilon_0} \frac{q}{r^2} \hat{r}
\end{equation*}

\subsection*{Campo Eléctrico debido a Múltiples Cargas Discretas}

\begin{equation*}
    \vec{E} = \sum_i \frac{1}{4\pi \epsilon_0} \frac{q_i}{r_i^2} \hat{r}_i
\end{equation*}

\subsection*{Campo Eléctrico debido a una Distribución Volumétrica de Carga}

\begin{equation*}
    \vec{E} = \frac{1}{4\pi \epsilon_0} \int \frac{\rho \, dV}{r^2} \hat{r}
\end{equation*}

\subsection*{Campo Eléctrico debido a una Línea de Carga}

\begin{equation*}
    \vec{E} = \frac{1}{4\pi \epsilon_0} \int \frac{\lambda \, dl}{r^2} \hat{r}
\end{equation*}

\subsection*{Campo Eléctrico debido a una Lámina de Carga}

\begin{equation*}
    \vec{E} = \frac{\sigma}{2\epsilon_0} \hat{n}
\end{equation*}

\section*{Densidad de Flujo Eléctrico, Ley de Gauss y Divergencia}

\subsection*{Densidad de Flujo Eléctrico}

\begin{equation*}
    \vec{D} = \epsilon_0 \vec{E}
\end{equation*}

\subsection*{Ley de Gauss}

\begin{equation*}
    \oint \vec{E} \cdot d\vec{A} = \frac{Q_{\text{int}}}{\epsilon_0}
\end{equation*}

\begin{equation*}
    \oint \vec{D} \cdot d\vec{A} = Q_{\text{int}}
\end{equation*}

\subsection*{Divergencia, Primera Ley de Maxwell}

\begin{equation*}
    \nabla \cdot \vec{E} = \frac{\rho}{\epsilon_0}
\end{equation*}

\begin{equation*}
    \nabla \cdot \vec{D} = \rho
\end{equation*}

\section*{Potencial Electrostático}

\begin{equation*}
    V = -\int \vec{E} \cdot d\vec{l}
\end{equation*}

\begin{equation*}
    \vec{E} = -\nabla V
\end{equation*}

\begin{equation*}
    \nabla^2 V = -\frac{\rho}{\epsilon_0}
\end{equation*}

\begin{equation*}
    \nabla^2 V = 0
\end{equation*}

\end{document}
